\documentclass[11pt]{article}

\usepackage[T1]{fontenc}
\usepackage{lmodern}
\usepackage{amsmath,amssymb,amsthm,mathtools}
\usepackage{array,booktabs,longtable}
\usepackage[margin=0.82in]{geometry}
\usepackage{microtype}
\usepackage[hidelinks]{hyperref}

\newtheorem{lemma}{Lemma}
\newtheorem{proposition}[lemma]{Proposition}
\theoremstyle{remark}
\newtheorem{remark}[lemma]{Remark}

\newcommand{\R}{\mathbb R}
\newcommand{\Weyl}{\mathcal W}
\newcommand{\qbar}{\overline q}
\newcommand{\pibar}{\overline\pi}
\newcommand{\piunder}{\underline\pi}

\title{Purely Analytic Ledger Replacement, Packet D:\\
Seam Arithmetic and the Two-Sided Payments}
\author{}
\date{Exact hand-checkable exhibit}

\begin{document}
\maketitle

\begin{abstract}
This packet replaces the twenty-two executable seam predicates and the
eleven moving/fixed payment predicates used in the high side through
$k_6$.  Every premise is an analytic inequality already derived in the
seam or staircase argument, and every remaining decision is displayed as
an exact rational identity or strict rational comparison.  In particular,
the derivative of the seam comparison function is derived explicitly and
is bounded by an exact proxy with a positive margin to $16/5$.  The six
moving $k_6$ payments are printed with their full rational differences,
not merely their numerators.  No program execution, floating-point value,
or interval calculation is a premise.
\end{abstract}

\section{Rational constants used by both tables}

For completeness, the needed bounds for $\pi$ can be checked from the
alternating arctangent series, rather than imported from an executable
Machin enclosure.  Machin's identity and the alternating remainder rule
give
\begin{align*}
 \pi
 &>16\left(\frac15-\frac1{3\,5^3}
       +\frac1{5\,5^5}-\frac1{7\,5^7}\right)-\frac4{239}
 =:L_\pi,\\
 \pi
 &<16\left(\frac15-\frac1{3\,5^3}
       +\frac1{5\,5^5}\right)
 -4\left(\frac1{239}-\frac1{3\,239^3}\right)
 =:U_\pi.
\end{align*}
Indeed, the tangent-addition formula proves
$\pi=16\arctan(1/5)-4\arctan(1/239)$, and the displayed truncations have
the indicated directions.  Direct rational subtraction gives
\begin{equation}
 L_\pi-\frac{333}{106}
 =\frac{3418213}{41563593750}>0,
 \qquad
 \frac{1571}{500}-U_\pi
 =\frac{323354809}{853244937500}>0.
 \label{eq:machin-gaps}
\end{equation}
Since
\begin{equation}
 \frac{22}{7}-\frac{1571}{500}=\frac3{3500}>0,
 \label{eq:pi-upper-chain}
\end{equation}
we may use throughout
\begin{equation}
 \piunder:=\frac{333}{106}<\pi<
 \pibar:=\frac{1571}{500}<\frac{22}{7}.
 \label{eq:pi-bounds}
\end{equation}
All later occurrences of a rational $\pi$ bound refer to
\eqref{eq:pi-bounds}.

\section{The twenty-two seam rows}

\subsection{Analytic input and comparison function}

We recall precisely the analytic data to which the finite arithmetic is
applied.  Put
\[
 \epsilon=1-\rho=y^2,\qquad
 \kappa=K\epsilon^2,\qquad
 0<y\leq\frac1{\sqrt6},\quad
 \rho\geq\frac56,\quad \kappa\geq54.
\]
For a first plateau of length $p$, and the quantities $m,R$ in the seam
decomposition, set
\[
 P=py,\qquad r=Ry,\qquad S=(p+m)y.
\]
The angular estimate $d>5/8$ and $R=p-dm$ imply, whenever $r>0$,
\begin{equation}
 0<r\leq P,\quad
 S\leq S_*(P,r):=\frac{13P-8r}{5}.
 \label{eq:Sstar}
\end{equation}
Equality in the second relation can occur when \(m=0\); no later step
requires strictness at this comparison.
The curvature-shelf calculation already proved in the seam argument gives
the self-consistent inequality
\begin{equation}
 P^2<\frac{2\pi^2Q}{(1-\widehat q)^2},\quad
 Q=1+\frac{y^2}{\kappa}+\frac{yS}{\kappa},\quad
 \widehat q=\frac{y^2}{\rho}(Q-1).
 \label{eq:self-consistent}
\end{equation}
Since the right-hand side increases with $Q$, define
\begin{equation}
 Q_*=1+\frac{y^2}{\kappa}+\frac{yS_*}{\kappa},\quad
 q_*=a(Q_*-1),\qquad a:=\frac{y^2}{\rho}\leq\frac15,
 \label{eq:Qstar}
\end{equation}
and
\begin{equation}
 H(P,r):=\frac{2\pi^2Q_*}{(1-q_*)^2}.
 \label{eq:H}
\end{equation}
Then \eqref{eq:Sstar} gives \(Q\leq Q_*\) and
\(\widehat q\leq q_*\).  Since \eqref{eq:self-consistent} is itself
strict and its right-hand side increases with \(Q\), it follows that
\(P^2<H(P,r)\).

The remaining arithmetic is exactly the following table.  Rows 1--5 are
the five constant/shelf rows, rows 6--10 are the five displacement rows,
rows 11--14 are the four derivative rows, rows 15--17 are the three
endpoint rows, and rows 18--22 are the five terminal-payment rows.  Thus
the table reconciles one-for-one with all twenty-two seam predicates.

For clarity, the radical used in rows 7--9 follows directly from the
analytic shelf bound.  Namely,
\[
 p^2<\frac{2\pi\rho K}{\epsilon},\quad P=py,\quad
 K=\frac{\kappa}{y^4}
\]
imply
\[
 \left(\frac{y^3P}{\kappa\rho}\right)^2
 <\frac{2\pi y^2}{\kappa\rho}
 \leq\frac{\pi}{135}<\frac{\pibar}{135}.
\]
Together with $S_*\leq13P/5$, this gives
\[
 q_*<\frac1{1620}+\frac{13}{5}\frac{763}{5000},
\]
which is precisely the dependency recorded by rows 6--9.

\begingroup
\small
\renewcommand{\arraystretch}{1.22}
\begin{longtable}{@{}>{\raggedleft\arraybackslash}p{0.035\linewidth}
 >{\raggedright\arraybackslash}p{0.24\linewidth}
 >{\raggedright\arraybackslash}p{0.47\linewidth}
 >{\raggedright\arraybackslash}p{0.15\linewidth}@{}}
\toprule
\# & quantity & exact rational check & role\\
\midrule
\endfirsthead
\toprule
\# & quantity & exact rational check & role\\
\midrule
\endhead
1 & seam $\pi$ bound &
$\pibar-U_\pi=323354809/853244937500>0$
& $\pi<\pibar$; \eqref{eq:machin-gaps}\\
2 & lower $\sqrt2$ bound &
$2-(140/99)^2=2/9801>0$
& $140/99<\sqrt2$\\
3 & upper $\sqrt2$ bound &
$(99/70)^2-2=1/4900>0$
& $\sqrt2<99/70$\\
4 & $y$ ceiling &
$(49/120)^2-1/6=1/14400>0$
& $y<49/120$\\
5 & shelf ceiling &
$(13/5)^2-2\pibar=119/250>0$
& $\sqrt{2\pi\rho}<13/5$; rows 1\\
\midrule
6 & first displacement term &
$\displaystyle\frac{(1/6)^2}{54(5/6)}=\frac1{1620}$
& $y^4/(\kappa\rho)\leq1/1620$\\
7 & radical ceiling &
$\displaystyle(763/5000)^2-\frac{\pibar}{135}
=\frac{8563}{675000000}>0$
& shelf radical; row 1\\
8 & displacement proxy &
$\displaystyle\frac1{1620}+\frac{13}{5}\frac{763}{5000}
=\frac{804689}{2025000}$
& rows 6--7\\
9 & displacement margin &
$\displaystyle\frac{159}{400}-\frac{804689}{2025000}
=\frac{497}{4050000}>0$
& $q_*,\widehat q<\qbar:=159/400$\\
10 & outer-half localization &
$\displaystyle\frac56(1-159/400)=\frac{241}{480}
=\frac12+\frac1{480}$
& $x_0/K>1/2$; row 9\\
\midrule
11 & $Q_*$-derivative proxy &
$\displaystyle D_Q:=\frac{13}{5}\frac{49/120}{54}
=\frac{637}{32400}$
& $\partial_PQ_*<D_Q$; row 4\\
12 & exact $H$-derivative proxy &
$\displaystyle
 \frac{2\pibar^2D_Q(1+159/400+2/5)}{(1-159/400)^3}
 =\frac{2260740364246}{708624500625}$
& rows 1, 9, 11 and $a\leq1/5$\\
13 & margin to $16/5$ &
$\displaystyle\frac{16}{5}
-\frac{2260740364246}{708624500625}
=\frac{6858037754}{708624500625}>0$
& $\partial_PH(P,B)<16/5$\\
14 & synthetic-slope margin &
$\displaystyle2\frac{14}{3}-\frac{16}{5}=\frac{92}{15}>0$
& $\partial_P(P^2-H)>92/15$\\
\midrule
15 & endpoint $Q_*$ proxy &
$\displaystyle
 \overline Q:=1+\frac{1/6}{54}
 +\frac{(49/120)(14/3)}{54}
 =\frac{10093}{9720},\quad
 \overline Q-1=\frac{373}{9720}$
& $Q_*(B,B)<\overline Q$\\
16 & endpoint $q_*$ proxy &
$\displaystyle\frac15(\overline Q-1)=\frac{373}{48600}$
& $q_*(B,B)<373/48600$\\
17 & initial reserve &
$\displaystyle
 \left(\frac{14}{3}\right)^2
 -\frac{2\pibar^2(10093/9720)}{(1-373/48600)^2}
 =\frac{2505132463469}{2616573970125}>0$
& $B^2-H(B,B)>0$; rows 1, 15--16\\
\midrule
18 & action-gain coefficient &
$\displaystyle
 g:=\frac{2(140/99)54}{3(1571/500)}
 =\frac{280000}{17281}$
& $K\eta>g/y$; rows 1--2\\
19 & gain after $R$ loss &
$\displaystyle g-\frac{14}{3}=\frac{598066}{51843}$
& use $R<14/(3y)$\\
20 & interface loss proxy &
$\displaystyle\frac8{15}\frac{99}{70}=\frac{132}{175}$
& $4\delta<132/175$; row 3\\
21 & dangerous-branch reserve &
$\displaystyle
 \frac{598066}{51843}\frac{120}{49}-1-\frac{132}{175}
 =\frac{80132733}{3024175}>0$
& $R>0$; rows 4, 19--20\\
22 & safe-branch reserve &
$\displaystyle
 \frac{280000}{17281}\frac{120}{49}-1-\frac{132}{175}
 =\frac{114694733}{3024175}>0$
& $R\leq0$; rows 4, 18, 20\\
\bottomrule
\end{longtable}
\endgroup

\subsection{Why the table proves the seam implications}

We spell out the two implications for which the table is more than a
collection of arithmetic facts.

First, on the segment $r=B:=14/3$ and $P\geq B$,
\[
 \partial_PQ_*=\frac{13y}{5\kappa},\quad
 \partial_Pq_*=a\,\partial_PQ_*.
\]
Differentiation of \eqref{eq:H}, followed only by the identity
$aQ_*=q_*+a$, gives
\begin{equation}
 \partial_PH
 =2\pi^2(\partial_PQ_*)
   \frac{1-q_*+2aQ_*}{(1-q_*)^3}
 =2\pi^2(\partial_PQ_*)
   \frac{1+q_*+2a}{(1-q_*)^3}.
 \label{eq:H-derivative}
\end{equation}
Rows 1, 9, and 11 and $a\leq1/5$ substitute monotonically into
\eqref{eq:H-derivative}; rows 12--13 then prove
$\partial_PH<16/5$.  Because $P\geq B$, row 14 proves
\begin{equation}
 \frac{d}{dP}\bigl(P^2-H(P,B)\bigr)
 =2P-\partial_PH(P,B)>\frac{92}{15}>0.
 \label{eq:synthetic-slope}
\end{equation}
Rows 15--17 prove $B^2-H(B,B)>0$.  Integration of
\eqref{eq:synthetic-slope} therefore yields $P^2>H(P,B)$ for every
$P\geq B$.  If $r\geq B$, then $H(P,r)\leq H(P,B)$ because
$S_*(P,r)$ decreases with $r$ and $H$ increases with $Q_*$.  This
contradicts $P^2<H(P,r)$.  Hence
\begin{equation}
 r<B,\quad R<\frac{14}{3y}.
 \label{eq:R-bound}
\end{equation}

Second, the turning-point action bound gives
\[
 K\eta\geq\frac{2\sqrt2\,\kappa}{3\pi}\frac1y
 >\frac{g}{y},
\]
where row 18 is the exact rational lower proxy.  Also row 4 gives
$1/y>120/49$, and row 20 gives
$4\delta<8\sqrt2/15<132/175$.  With
$M=\lfloor K\eta\rfloor-R$ and the strict inequality
$\lfloor x\rfloor>x-1$, \eqref{eq:R-bound} and row 21 give, when
$R>0$,
\[
 M-4\delta>
 \left(g-\frac{14}{3}\right)\frac{120}{49}
 -1-\frac{132}{175}>0.
\]
When $R\leq0$, row 22 gives the still stronger estimate
\[
 M-4\delta>
 g\frac{120}{49}-1-\frac{132}{175}>0.
\]
Thus the seam tail has the strict action payment claimed in the main
argument, including the $R=0$ and integer-$K\eta$ faces.

\begin{remark}[Notation correction in the executable label]
The supplemental verifier labels row 11 as an ``exact $dq/dP$ cap.''
In the notation of the proof, the number $637/32400$ is a proxy for
$\partial_PQ_*$, not for $\partial_Pq_*$.  In fact
$\partial_Pq_*=a\partial_PQ_*$ and hence
$\partial_Pq_*<637/162000$.  Formula \eqref{eq:H-derivative} and the
verifier both use $637/32400$ in its correct role as the $Q_*$-derivative
proxy.  The proof is therefore sound, but the old ledger label is
inaccurate.
\end{remark}

\begin{remark}[Source wording for the $\pi$ bound]
One sentence in the seam passage groups $\pi<1571/500$ with bounds said
to be proved by positive rational squaring.  Squaring proves rows 2--5,
but not a bound for $\pi$.  Equations \eqref{eq:machin-gaps}--
\eqref{eq:pi-bounds} supply the required analytic alternating-series
proof.
\end{remark}

\section{The eleven \texorpdfstring{$k_6$}{k6} payment rows}

Put
\begin{equation}
 z_\rho=\frac{\pi}{1-\rho},\quad
 k_m(\rho)^2=z_\rho^2+m(m+1),\quad
 p_1(\rho)^2=4z_\rho^2+2,
 \label{eq:walls}
\end{equation}
and
\begin{equation}
 \Weyl(\rho,K)=\frac{2}{9\pi}(1-\rho^3)K^3.
 \label{eq:W}
\end{equation}
For a moving wall $T^2=az_\rho^2+b$, direct differentiation shows that
$\Weyl(\rho,T)$ increases whenever
\begin{equation}
 a\pi^2(1+\rho)>b\rho^2(1-\rho)^2.
 \label{eq:moving-sign}
\end{equation}
Here $b/a\leq42$.  Since
$\rho^2(1-\rho)^2\leq1/16$ and $\pi>3$, the left side of
\eqref{eq:moving-sign}, after division by $a$, is greater than $9$,
whereas the right side is at most $42/16$; hence every moving payment in
the table is increasing.  At a fixed wall $L$, the function
$\Weyl(\rho,L)$ is decreasing in $\rho$.

For a rational split $r$, define
\begin{equation}
 A(a,b;r):=a\frac{(333/106)^2}{(1-r)^2}+b,\quad
 B(r):=\frac7{99}(1-r^3).
 \label{eq:AB}
\end{equation}
The bounds \eqref{eq:pi-bounds} give, on a moving wall at $\rho=r$,
\begin{equation}
 \Weyl(r,T)>B(r)A(a,b;r)^{3/2}.
 \label{eq:moving-lower}
\end{equation}
Thus, for a positive cap $C$, the rational inequality
\begin{equation}
 \Delta_{\rm mov}:=B(r)^2A(a,b;r)^3-C^2>0
 \label{eq:moving-reserve}
\end{equation}
implies the strict payment, because
\[
 \Delta_{\rm mov}
 =\bigl(B(r)A^{3/2}-C\bigr)
  \bigl(B(r)A^{3/2}+C\bigr)
\]
and the second factor is positive.  At a fixed wall $L$, it is enough to
check the ordinary rational reserve
\begin{equation}
 \Delta_{\rm fix}:=B(r)L^3-C>0.
 \label{eq:fixed-reserve}
\end{equation}

The following is the complete eleven-row registry: six moving rows and
five fixed rows.  The fractions in the last column are the full reduced
differences in \eqref{eq:moving-reserve} or \eqref{eq:fixed-reserve}.

\begingroup
\scriptsize
\renewcommand{\arraystretch}{1.25}
\begin{longtable}{@{}>{\raggedright\arraybackslash}p{0.09\linewidth}
 >{\centering\arraybackslash}p{0.09\linewidth}
 >{\centering\arraybackslash}p{0.115\linewidth}
 >{\centering\arraybackslash}p{0.07\linewidth}
 >{\raggedright\arraybackslash}p{0.53\linewidth}@{}}
\toprule
wall/branch & split $r$ & wall data & cap $C$ & full exact rational difference\\
\midrule
\endfirsthead
\toprule
wall/branch & split $r$ & wall data & cap $C$ & full exact rational difference\\
\midrule
\endhead
$k_1$ moving & $1/8$ & $(a,b)=(1,2)$ & $4$ &
$\displaystyle\Delta_{\rm mov}=
\frac{153864824754340538785}{348796101192617852928}>0$\\
$k_2$ moving & $3/10$ & $(a,b)=(1,6)$ & $9$ &
$\displaystyle\Delta_{\rm mov}=
\frac{1399810848073457853053}{394237491401523000000}>0$\\
$k_3$ moving & $1/2$ & $(a,b)=(1,12)$ & $16$ &
$\displaystyle\Delta_{\rm mov}=
\frac{137027505353736631}{514922437748928}>0$\\
$k_4$ moving & $1/2$ & $(a,b)=(1,20)$ & $25$ &
$\displaystyle\Delta_{\rm mov}=
\frac{2507119282010251109}{13902905819221056}>0$\\
$k_5$ moving & $13/25$ & $(a,b)=(1,30)$ & $36$ &
$\displaystyle\Delta_{\rm mov}=
\frac{92672404686738742434741}{709259556128000000000}>0$\\
$p_1$ moving & $1/5$ & $(a,b)=(4,2)$ & $29$ &
$\displaystyle\Delta_{\rm mov}=
\frac{373255772037478993002833}{868931613701316000000}>0$\\
\midrule
$k_2$ fixed & $3/10$ & $L=51/10$ & $9$ &
$\displaystyle\Delta_{\rm fix}=\frac{1387329}{11000000}>0$\\
$k_3$ fixed & $1/2$ & $L=13/2$ & $16$ &
$\displaystyle\Delta_{\rm fix}=\frac{6277}{6336}>0$\\
$k_4$ fixed & $1/2$ & $L=15/2$ & $25$ &
$\displaystyle\Delta_{\rm fix}=\frac{775}{704}>0$\\
$k_5$ fixed & $13/25$ & $L=17/2$ & $36$ &
$\displaystyle\Delta_{\rm fix}=\frac{452843}{343750}>0$\\
$p_1$ fixed & $1/5$ & $L=77/10$ & $29$ &
$\displaystyle\Delta_{\rm fix}=\frac{849901}{281250}>0$\\
\bottomrule
\end{longtable}
\endgroup

We now verify that the table pays every branch rather than merely listing
positive numbers.  For each maximum entry wall
\[
 \max\{k_2,51/10\},\quad
 \max\{k_3,13/2\},\quad
 \max\{k_4,15/2\},\quad
 \max\{k_5,17/2\},\quad
 \max\{p_1,77/10\},
\]
partition the ratio interval at the displayed $r$.  On $\rho\geq r$,
the maximum is at least the moving component and monotonicity together
with the moving row pays $C$.  On $\rho\leq r$, the maximum is at least
the fixed component and the fixed row pays $C$.  This argument does not
assume that the two components are equal at the rational split.  For
$k_1$, the applicable ratio interval starts at
$\rho_c>1/8$ (because $\pi<22/7<7/2$), so the first moving row applies
directly.  Strict payments also cover coincident entry walls, while the
wall itself remains assigned to the lower cap by the strict spectral
convention.

\section{Reconciliation with the lower-\texorpdfstring{$d$}{d} rows}

The two rational $\pi$ predicates used by the two-sided ledger are already
proved in \eqref{eq:pi-bounds}.  The eleven $k_6$ predicates are exactly
the eleven rows in the preceding table.  The remaining thirteen
wall/cap predicates in that ledger are the ten transparent multiplicity
sums
\begin{equation}
\begin{aligned}
 &(1,\ 1+3,\ 1+3+5,\ 1+3+5+1,\ 1+3+5+7,\\
 &\quad 1+3+5+7+1,\ 1+3+5+7+9+1,\\
 &\quad 1+3+5+7+9+1+3,\\
 &\quad 1+3+5+7+9+11+1+3,\\
 &\quad 1+3+5+7+9+11+1+3+5)\\
 &\hspace{35mm}=(1,4,9,10,16,17,26,29,40,45),
\end{aligned}
 \label{eq:lower-caps}
\end{equation}
and the endpoint ordering $13/2<2z_\rho<15/2$.  The latter follows
without an executable check: using $\sqrt{337}<19$ on the lower endpoint
and \eqref{eq:pi-bounds},
\begin{equation}
 2\frac{333}{106}\frac{19}{18}-\frac{13}{2}
 =\frac7{53}>0,
 \qquad
 \frac{15}{2}-\left(2\frac{22}{7}+1\right)
 =\frac3{14}>0.
 \label{eq:2z-order}
\end{equation}
Thus the cap--$10$ row is empty, and the three executable ordering rows
collapse to \eqref{eq:2z-order} and its immediate empty-cell consequence.

The nine lower-$d$ payment rows need no new table here: the main paper's
lemma ``Lower side through $d$'' (label
\texttt{cm:lem-lower-d}) already prints the seven fixed-wall differences
\[
 \frac{793292}{1546875},\ 
 \frac{486236661}{1375000000},\ 
 \frac{333100003}{99000000},\ 
 \frac{329797}{88000},\ 
 \frac{3590476297}{1125000000},\ 
 \frac{327078887}{99000000},\ 
 \frac{8805519}{1375000},
\]
and states the initial and moving payments
$\Weyl(\rho_c,L(\rho_c))>19/6$, where
$L(\rho)=(2\rho)^{-1}$ on this branch, and
$\Weyl(1/19,2z)>508/27$.  If desired, their exact printed lower proxies close
by the two one-line subtractions
\begin{equation}
 \frac{57421}{16854}-\frac{19}{6}=\frac{675}{2809}>0,
 \qquad
 \frac{173863}{8427}-\frac{508}{27}
 =\frac{137795}{75843}>0.
 \label{eq:lower-d-last-two}
\end{equation}
Consequently every seam and two-sided payment predicate assigned to Packet
D is now visible as analytic reasoning plus printed integer arithmetic.

\section{Dependency count and conclusion}

The replacement count is exact:
\[
 \underbrace{22}_{\text{seam table}}
 +\underbrace{(2+6+5+10+3)}_{k_6/\text{lower-}d\text{ wall--cap ledger}}
 +\underbrace{9}_{\text{already printed lower-}d\text{ payments}}.
\]
The two constant rows are supplied by \eqref{eq:pi-bounds}; the six plus
five payment rows are the eleven-row table; the ten bookkeeping rows are
\eqref{eq:lower-caps}; the three ordering rows are
\eqref{eq:2z-order}; and the nine already printed payments are explicitly
cross-referenced above.  This removes program execution from the logical
dependency of precisely the Packet-D portion of the exact ledger.

\end{document}
